%% Skeleton outline of the paper, in the style of references/outline.pdf.
%% Build: latexmk -pdf skeleton.tex
\documentclass[11pt,a4paper]{article}

\usepackage[T1]{fontenc}
\usepackage{libertine}
\usepackage[varqu]{zi4}
\usepackage[margin=2.4cm]{geometry}
\usepackage{xcolor}
\usepackage{enumitem}
\usepackage[hidelinks]{hyperref}

\setlist{nosep,leftmargin=2em}
\setlength{\parindent}{0pt}
\setlength{\parskip}{0.3em}

\newcommand{\tool}{\texttt{spec4pbt}}
\newcommand{\desc}[1]{{\small #1\par}}
\newcommand{\fig}[2]{{\small\color{gray!70!black}\textsf{Fig.~#1}\quad #2\par}}

\title{Toward Trustworthy Agentic Property-Based Testing of\\
  Language Implementations from Mechanized Specifications}
\date{}

\begin{document}
\maketitle
\vspace{-3em}

\section{Introduction}
\textbf{Motivation:}
\begin{enumerate}\small
  \item PBT needs properties and generators; agents can now build generators from user-supplied properties.
  \item Trusting the resulting tests requires a valid property and a correct generator.
  \item Mechanized specifications let us state and prove both, and an ITP checks the proofs before any test runs.
\end{enumerate}
\textbf{Contributions:}
\begin{enumerate}\small
  \item three challenges for trustworthy PBT of language implementations from mechanized specifications (\S2.2);
  \item \tool{}, which addresses them for SpecTec specifications with Lean (\S3);
  \item a case study on P4, compared with Failing with Purpose (\S5).
\end{enumerate}

\section{Motivating Example}
\subsection{Example: Generating Tests for P4 Implementations}
\desc{Well-typed acceptance and ill-typed rejection are example properties of P4 implementations within a stated language fragment and target.}
\desc{P4Smith generates valid P4 programs; Failing with Purpose generates ill-typed ones by mutation. Both examples raise questions about property validity, the match between possible generator outputs and the property's precondition, and the interpretation of implementation observations.}

\subsection{Three Challenges for Trustworthy PBT}
\begin{enumerate}\small
  \item \textbf{Properties:} state the property precisely and prove it over the specification.
  \item \textbf{Generators:} prove the generator sound and complete for the property's precondition; assess efficiency and sampling diversity separately.
  \item \textbf{Implementations:} derive what to run and what to check from the property's goal.
\end{enumerate}

\section{\tool: Trustworthy PBT Toolchain for SpecTec Specifications}
\subsection{Overview}
\desc{The agent constructs proofs and generators, and Lean checks them.}
\fig{1}{Overview: orange arrows are agent construction, blue arrows are Lean checking.}

\subsection{SpecTec and Its Translation to Lean}
\desc{Syntax becomes inductive types, and rules become clauses tried in order.}
\fig{2}{P4 \texttt{Cast\_impl\_neq} in SpecTec and its ordered Lean translation.}

\subsection{Formalizing and Proving Properties}
\desc{A SpecTec-like property yields a Lean proposition and an executable checker.}
\fig{3}{Instantiation property and its derived Lean contract.}

\subsection{Constructing and Verifying Generators}
\desc{Following Palamedes, the agent builds the generator by proving that a sound and complete one exists.}

\subsection{Testing External Implementations}
\desc{User-supplied adapters map implementation outcomes to relation results, and the goal turns them into verdicts.}

\section{Towards Decomposition}
\subsection{Problem}
\desc{Constructing and verifying a generator as a single task is difficult because specification constraints interact.}
\desc{Soundness and completeness alone do not rule out inefficient generate-then-check strategies.}
\subsection{Analyzing Generation Order}
\desc{Generation efficiency depends on both variable order and the choice of inputs and outputs for each relation.}
\desc{We propose algorithmically analyzing variable dependencies to determine what to generate first and which constraints each step should establish.}
\subsection{Towards Agent-Guided Composition}
\desc{Order analysis alone cannot enforce correct-by-construction generation.}
\desc{We therefore explore using the analysis to guide an agent in constructing and composing smaller generators, while Lean checks the original soundness and completeness obligations.}
\desc{The concrete decomposition procedure remains open.}

\section{Evaluation}
\desc{RQ1--RQ3 study ill-typed rejection; RQ4 uses several properties.}

\subsection{Setup and Validation}
\desc{Targets, agent settings, and budgets; the translation is validated on \texttt{p4c} tests.}

\subsection{RQ1: Feasibility}
\desc{Can an agent prove the property and build a verified generator, with and without decomposition?}

\subsection{RQ2: Hallucinations Caught by Lean}
\desc{Which wrong proofs and generators would Lean's checks have kept out of testing?}

\subsection{RQ3: Generator Efficiency}
\desc{Throughput and coverage against Failing with Purpose, Palamedes, and QuickChick.}

\subsection{RQ4: Bug Finding}
\desc{Bugs found with several properties, including ones on recent P4 features.}

\section{Discussion}
\subsection{Expressiveness of Properties}
\subsection{Guarantees and Limitations}
\subsection{Generalization Beyond P4}
\desc{Other SpecTec languages (Wasm, Ethereum consensus) and other specification languages.}

\section{Related Work}
\begin{enumerate}\small
  \item Testing from mechanized specifications (Redex, WEST, Failing with Purpose) does not prove the property or the generator.
  \item PBT in interactive theorem provers (QuickChick, Palamedes) tests conjectures before proving them.
\end{enumerate}

\section{Conclusion}

\end{document}
